\section{问题四：实时波动电价下的购电策略}

\subsection{问题分析}

问题四要求在外网电价实时波动的条件下，重新制定问题二与问题三中的购电与储能调度策略。与前几问不同，当天真实电价在制定购电计划时尚未完全获得，而实际费用又需按照实时电价结算，因此本问首先需要解决未来电价未知条件下的价格预测问题，并在此基础上重新进行购电优化。

考虑到附件~4 中的电价具有较稳定的日内变化规律，同时受到星期和近期价格水平的影响，本文根据决策日前已经获得的历史电价，采用历史同刻均值、星期修正和近期修正相结合的方法预测当天各时段电价，以保证预测过程不使用未来信息。随后分别沿用问题二和问题三的模型，根据预测电价以及负载、光伏预测结果制定购电与储能调度计划。最后，利用附件~4 的实际电价计算全年实际运行费用，通过总费用、紧急购电量及不同策略间的费用差异评价购电策略。

\subsection{日前电价预测模型建立}

问题四中，每日 $0{:}00$ 制定购电计划时尚无法获得当天真实电价，因此需要利用此前的历史电价预测当天 $144$ 个 $10\,\mathrm{min}$ 时段的价格。考虑到附件~4 的实时电价具有较稳定的日内规律，同时存在星期差异和近期波动，本文将预测电价分解为历史同刻基础价格、星期修正和近期修正三部分。

记第 $d$ 日第 $t$ 个时段的真实电价和预测电价分别为 $p_{d,t}$ 和 $\hat p_{d,t}$，则
\begin{equation}
\hat p_{d,t}
=
b_{d,t}
+
\Delta_{\mathrm{week},d,t}
+
\Delta_{\mathrm{recent},d,t},
\qquad t=1,\ldots,144.
\label{eq:q4_price_forecast}
\end{equation}

其中，$b_{d,t}$ 为历史同刻基础价格，$\Delta_{\mathrm{week},d,t}$ 和 $\Delta_{\mathrm{recent},d,t}$ 分别为星期修正项和近期修正项。

\textbf{（1）历史同刻基础价格}

为提取电价稳定的日内变化规律，对决策日前所有历史日期在相同时段的真实电价取平均。设第 $d$ 日以前共有 $n_d$ 个历史日期，则
\begin{equation}
b_{d,t}
=
\frac{1}{n_d}
\sum_{j<d}p_{j,t}.
\label{eq:q4_price_base}
\end{equation}

\textbf{（2）星期修正项}

为进一步反映不同星期日期之间的价格差异，定义第一层残差
\begin{equation}
r_{j,t}
=
p_{j,t}-b_{j,t}.
\label{eq:q4_price_residual_week}
\end{equation}

记 $w(d)$ 为第 $d$ 日的星期编号，则取决策日前相同星期日期的第一层残差均值作为星期修正：
\begin{equation}
\Delta_{\mathrm{week},d,t}
=
\frac{1}{|\mathcal W_d|}
\sum_{j\in\mathcal W_d}r_{j,t},
\qquad
\mathcal W_d=\{j<d\mid w(j)=w(d)\}.
\label{eq:q4_week_correction}
\end{equation}

\textbf{（3）近期修正项}

在扣除历史同刻价格和星期影响后，定义第二层残差
\begin{equation}
e_{j,t}
=
p_{j,t}
-
b_{j,t}
-
\Delta_{\mathrm{week},j,t}.
\label{eq:q4_price_residual_recent}
\end{equation}

取决策日前最近 $7$ 个已观测日期的第二层残差进行等权平均，得到近期修正项：
\begin{equation}
\Delta_{\mathrm{recent},d,t}
=
\frac{1}{m_d}
\sum_{k=1}^{m_d}e_{d-k,t},
\qquad
m_d=\min(7,n_d).
\label{eq:q4_recent_correction}
\end{equation}

该修正用于反映近期价格变化，并通过多日平均减弱单日异常波动的影响。

由此，模型依次提取长期日内规律、星期差异和近期价格变化。

\textbf{（4）预测模型预热与滚动更新}

附件~4 的实时电价从 $2025$ 年 $1$ 月 $1$ 日开始记录，而正式调度从 $2$ 月 $1$ 日开始，因此将 $1$ 月作为预测模型的预热区间。由于 $1$ 月 $1$ 日缺少更早的历史电价，采用附件~1 的分时电价 $p_t^{(1)}$ 进行初始化：
\begin{equation}
b_{1,t}=p_t^{(1)},\qquad
\Delta_{\mathrm{week},1,t}=0,\qquad
\Delta_{\mathrm{recent},1,t}=0.
\label{eq:q4_price_cold_start}
\end{equation}

此后按日期顺序逐日更新。每天利用截至前一日的历史信息计算预测电价，并在当天真实电价获得后更新历史价格及残差。到 $2$ 月 $1$ 日正式调度开始时，基础价格由 $1$ 月 $31$ 天的同刻真实电价计算，星期修正和近期修正则利用 $1$ 月逐日形成的历史残差计算。1 月仅用于模型预热，不计入正式费用统计。

正式调度期间继续采用相同的逐日更新方式，每日 $0{:}00$ 得到预测电价向量
\begin{equation}
\hat{\boldsymbol p}_d
=
\left(
\hat p_{d,1},
\hat p_{d,2},
\ldots,
\hat p_{d,144}
\right)^{\mathsf T}.
\label{eq:q4_price_vector}
\end{equation}

预测电价用于当天购电决策，附件~4 的真实电价仅用于实际费用计算，从而避免在决策过程中使用未来电价信息。

\textbf{（4）初始历史数据与逐日更新}

附件~4 的实时电价从 $2025$ 年 $1$ 月 $1$ 日开始记录，而正式调度从 $2$ 月 $1$ 日开始，因此利用 $1$ 月数据积累电价预测所需的历史信息。由于 $1$ 月 $1$ 日缺少更早的历史电价，采用附件~1 的分时电价 $p_t^{(1)}$ 进行初始化：
\begin{equation}
b_{1,t}=p_t^{(1)},\qquad
\Delta_{\mathrm{week},1,t}=0,\qquad
\Delta_{\mathrm{recent},1,t}=0.
\label{eq:q4_price_cold_start}
\end{equation}

此后按日期顺序逐日计算预测电价，并在当天真实电价获得后更新历史价格及残差。到 $2$ 月 $1$ 日时，基础价格由 $1$ 月 $31$ 天的同刻真实电价计算，星期修正和近期修正则根据 $1$ 月逐日形成的历史残差计算。1 月仅用于积累上述历史信息，不计入正式费用统计。

正式调度期间，每日继续利用截至前一日的历史信息更新预测，得到当天的预测电价向量
\begin{equation}
\hat{\boldsymbol p}_d
=
\left(
\hat p_{d,1},
\hat p_{d,2},
\ldots,
\hat p_{d,144}
\right)^{\mathsf T}.
\label{eq:q4_price_vector}
\end{equation}

预测电价用于制定当天购电计划，附件~4 的真实电价用于实际费用计算，以保证决策时不使用尚未获得的电价信息。


\subsection{模型求解}

问题四的正式费用统计区间为 $2025$ 年 $2$ 月 $1$ 日至 $12$ 月 $31$ 日，共 $334$ 天。求解过程按照“电价预测、计划优化、实际回放、费用结算和状态更新”的顺序逐日进行。问题四沿用问题二和问题三已经确定的负荷预测、光伏预测和储能控制方法，将原模型中的固定分时电价替换为因果预测电价，并采用附件~4 的真实电价进行最终结算。

\textbf{Step 1：生成当日电价预测。}

首先读取附件~4 的全年实时电价，并按照日期顺序模拟每日 $0{:}00$ 的因果预测。对于决策日 $d$，模型只使用日期早于 $d$ 的历史电价，依次计算历史同刻基础价格、星期修正和最近 $7$ 日修正，得到当天的预测电价向量
\begin{equation}
\hat{\boldsymbol p}_{d}
=
\left(
\hat p_{d,1},
\hat p_{d,2},
\ldots,
\hat p_{d,144}
\right)^{\mathsf T}.
\end{equation}

在 $1$ 月 $1$ 日缺少更早历史数据时，采用附件~1 的分时电价曲线完成冷启动；此后每日先生成预测，再将当天已经实现的附件~4 电价及由此计算的两层残差加入历史统计。到 $2$ 月 $1$ 日时，基础价格已经使用 $1$ 月 $1$ 日至 $1$ 月 $31$ 日的同刻真实电价计算，星期修正和近期修正则由 $1$ 月逐日形成的历史残差确定。1 月电价预测不计入正式费用，但其真实价格和因果残差用于正式区间的电价预测。

\textbf{Step 2：求解对应问题二的日前调度模型。}

对于问题四中对应问题二的部分，$2$ 月 $1$ 日 $0{:}00$ 的初始储电量设为
\begin{equation}
E_{\mathrm{Feb.1},0}=6000\,\mathrm{kWh}.
\end{equation}
此后各日的初始储电量均取上一日实际执行结束时的储电量，从而保持储能状态的跨日连续性。

每日 $0{:}00$，采用问题二的因果扩展窗 XGBoost 模型预测当天负荷，并以过去 $7$ 日同一时段光伏功率的均值预测当天光伏出力。在此基础上，对历史净负荷残差取 $\alpha=0.8$ 分位数，形成用于优化的保守净负荷。以预测电价 $\hat p_{d,t}$ 为目标函数系数，求解日前线性规划
\begin{equation}
\min C_{d}^{\mathrm{plan}}
=
\sum_{t=1}^{144}\hat p_{d,t}G_{d,t},
\end{equation}
并在电量平衡、储能容量、接口功率及非负约束下，联合确定全天计划购电量和计划储能轨迹。当前正式策略还设置计划日末储电量硬约束
\begin{equation}
E_{d,144}^{\mathrm{plan}}
=S^{*}=2400\,\mathrm{kWh}.
\end{equation}

求解得到的计划购电量 $G_{d,t}$ 在当天保持不变。计划充放电量及计划储电量仅描述预测场景下的最优轨迹，实际充放电量仍需在回放阶段根据当天真实负荷和光伏逐时段确定。

\textbf{Step 3：回放对应问题二的实际运行过程。}

锁定全天计划购电量后，按照时间顺序代入当天真实负荷和光伏数据，逐时段执行储能充放电。当计划购电与光伏发电形成电量盈余时，优先在储能容量和接口功率允许的范围内充电，不能存入储能的剩余电量作为弃电；当二者不足以满足负荷时，按照 $\rho=0.625$ 的计划轨迹保护规则放电。具体而言，实际储电量不得低于由储能下限与计划储电量共同确定的动态保护水平
\begin{equation}
E_{d,t}^{\mathrm{res}}
=E_{\min}+\rho\left(E_{d,t}^{\mathrm{plan}}-E_{\min}\right),
\qquad \rho=0.625,
\end{equation}
从而避免储能在较早时段被过度消耗。

若计划购电、实际光伏和允许的储能放电仍不足以满足负荷，则不足部分记为紧急购电量 $e_{d,t}^{\mathrm{em}}$。当天实际费用按照附件~4 的真实电价计算：
\begin{equation}
C_{d}^{(4-2)}
=
\sum_{t=1}^{144}\pi_{d,t}G_{d,t}
+
5\sum_{t=1}^{144}\pi_{d,t}e_{d,t}^{\mathrm{em}},
\label{eq:q42_actual_cost}
\end{equation}
其中，$\pi_{d,t}$ 为附件~4 中第 $d$ 日第 $t$ 个时段的真实电价。实际日末储电量作为下一日的初始储电量，形成跨日连续的实际运行轨迹。

\textbf{Step 4：求解对应问题三的日内滚动调度模型。}

对于问题四中对应问题三的部分，沿用问题三的正式 LA 策略。该策略不进行 $1$ 月储能运行预热，而是与 Q4-2 采用相同的正式起点，直接设置
\begin{equation}
E_{\mathrm{Feb.1},0}=6000\,\mathrm{kWh}.
\end{equation}
1 月数据只用于电价预测以及负荷、光伏预测残差分位数的历史积累，不参与 Q4-3 的储能状态传递和费用统计。

每日 $0{:}00$，采用因果扩展窗 XGBoost 模型生成当天负荷点预测，并使用附件~3 在 $0{:}00$ 发布的光伏预报形成初始购电合同。随后分别在 $6{:}00$、$12{:}00$ 和 $18{:}00$ 进行日内更新：根据当前时刻以前实际负荷与 $0{:}00$ 负荷预测的累计均值之比，对当天剩余时段的负荷预测作比例修正；同时读取附件~3 在相应时刻发布的最新光伏预报，并以发布前最后一个已经完成的 $10\,\mathrm{min}$ 时段实测光伏功率作为插值起点，将小时级预报转换为 $10\,\mathrm{min}$ 序列。

LA 策略采用最大 $48\,\mathrm{h}$ 展望。在 $0{:}00$，优化范围包括当天和次日，共 $288$ 个 $10\,\mathrm{min}$ 时段；在 $6{:}00$、$12{:}00$ 和 $18{:}00$，优化范围分别由当前时刻延伸至次日 $24{:}00$。由于次日尚未发布对应的负荷和光伏预报，次日展望部分采用目标日期前一周同一星期日期的实际负荷与光伏曲线构造因果预测。

对于每次滚动优化，下一次更新时间之前的 $6$ 小时锁定段采用净负荷残差分位数 $q_L=0.60$；当天以后仍可调整的时段及次日展望时段采用 $q=0.50$。各分位修正均由决策日以前已经闭合的预测残差计算。当天四次优化始终使用 $0{:}00$ 得到的预测电价 $\hat{\boldsymbol p}_{d}$，不将附件~4 中当天已经发生或未来时段的真实电价作为优化输入；对于次日展望部分，则将当天的 $144$ 段预测价格按照相同日内时刻循环扩展。各次调整只改变尚未执行时段的购电合同，已经执行的时段保持不变。

在相邻两次更新之间，将当前已经确定的合同购电量与真实负荷、光伏逐时段回放，并在容量及接口功率约束下贪婪执行储能充放电。当天实际日末储电量继续传递至下一日。

\textbf{Step 5：按照真实电价进行合同偏差结算。}

记每日 $0{:}00$ 确定的初始计划购电量为 $G_{d,t}$，经过日内滚动调整后最终执行的合同购电量为 $G'_{d,t}$，定义增购量和减购量分别为
\begin{equation}
\Delta_{d,t}^{+}
=
\max\left(G'_{d,t}-G_{d,t},0\right),
\qquad
\Delta_{d,t}^{-}
=
\max\left(G_{d,t}-G'_{d,t},0\right).
\end{equation}

初始计划高于最终合同的部分按照真实电价的 $0.5$ 倍退回，最终合同高于初始计划的部分按照真实电价的 $1.5$ 倍增购；实际运行中仍存在的供电缺口按照真实电价的 $5$ 倍紧急购买。因此，问题四对应问题三的每日实际费用为
\begin{equation}
\begin{aligned}
C_{d}^{(4-3)}
={}&
\sum_{t=1}^{144}\pi_{d,t}G_{d,t}
-
0.5\sum_{t=1}^{144}\pi_{d,t}\Delta_{d,t}^{-}\\
&+
1.5\sum_{t=1}^{144}\pi_{d,t}\Delta_{d,t}^{+}
+
5\sum_{t=1}^{144}\pi_{d,t}e_{d,t}^{\mathrm{em}}.
\end{aligned}
\label{eq:q43_actual_cost}
\end{equation}

该结算只在调度方案确定后进行。预测电价决定购电与储能调度，真实电价只用于计算最终费用，从而避免在决策阶段引入当天尚未获得的价格信息。

\textbf{Step 6：检查求解状态并汇总结果。}

上述日前规划与日内滚动规划均为连续线性规划，采用 CBC 求解器求解，单次求解时间上限为 $60\,\mathrm{s}$。只有当求解器返回状态为 \texttt{Optimal} 时，才提取购电量、充放电量和储电量轨迹并继续实际运行回放；当前正式结果中的各次线性规划均获得最优状态。

完成逐日求解后，将所得方案重新代入电量平衡、储电量状态转移、储能容量、接口功率及非负约束，检查实际 SOC 的跨日连续性，并按照附件~4 的真实电价独立复算各项费用。同时逐日核验电价预测和负荷、光伏预测均未使用决策时点以后才能获得的信息。当前 Q4-2 的约束校验错误数和价格泄漏数均为 $0$，Q4-3 的约束门禁错误与信息泄漏错误也均为空。最后汇总 $334$ 个正式日期的计划购电费用、合同增减费用、紧急购电费用、紧急购电量及全年总费用，用于后续策略比较和结果分析。

\subsection{求解结果}

通过上述预测与调度过程，分别得到实时波动电价下问题二的日前购电方案和问题三的日内滚动购电方案。表~\ref{tab:q4_all} 给出了回测区间内两类策略及其对照方案的全年费用与紧急购电情况。

\begin{table}[!htbp]
  \centering
  \caption{波动电价下问题二与问题三的全年结果比较}
  \label{tab:q4_all}
  \small

  \textbf{(a) Q4-2：日前购电策略}\par\medskip
  \begin{tabular}{lrrr}
    \toprule
    方案 & 总费用（元） & 紧急购电量（kWh） & 节省费用（元） \\
    \midrule
    C\_pv7d\_q82（固定 $q=0.8$）
      & $14\,480\,739.96$ & $65\,421.48$ & --- \\
    \textbf{E\_pv7d\_netrho（本文）}
      & \textbf{$14\,413\,479.02$} & $85\,500.89$ & \textbf{$67\,260.94$} \\
    \bottomrule
  \end{tabular}

  \medskip

  \textbf{(b) Q4-3：日内滚动购电策略}\par\medskip
  \begin{tabular}{lrrr}
    \toprule
    方案 & 总费用（元） & 紧急购电量（kWh） & 较 N0 节省（元） \\
    \midrule
    N0（不调整）
      & $15\,983\,608.51$ & $466\,576.26$ & --- \\
    M0L（选择门）
      & $14\,826\,830.64$ & $218\,215.59$ & $1\,156\,777.87$ \\
    \textbf{LA（本文）}
      & \textbf{$14\,001\,199.26$} & \textbf{$67\,153.26$} & \textbf{$1\,982\,409.26$} \\
    \bottomrule
  \end{tabular}

  \par\medskip
  {\small 注：Q4-3 本文 LA 的合同购电费用为 $12\,979\,047.00$ 元，紧急购电费用为 $370\,783.78$ 元，其余费用来自少买与多买偏差结算。各方案均采用附件~4 的真实电价按相同规则回测。}
\end{table}

由表~\ref{tab:q4_all}(a) 可见，Q4-2 本文方案全年总费用为 $14\,413\,479.02$ 元，较固定 $q=0.8$ 方案减少 $67\,260.94$ 元，降幅约为 $0.46\%$。与此同时，紧急购电量由 $65\,421.48$ kWh 增加至 $85\,500.89$ kWh，说明本文策略并非单纯追求减少紧急购电，而是在计划购电成本与高价紧急购电风险之间进行权衡，最终取得更低的全年总费用。

图~\ref{fig:q4_dispatch} 进一步给出了 Q4-2 本文方案的日级运行情况。

\Needspace{10\baselineskip}
\begin{figure}[!htbp]
  \centering
  \includegraphics[width=0.72\textwidth]{q4_official_dispatch_details.png}
  \caption{Q4-2 本文方案日级运行情况}
  \label{fig:q4_dispatch}
\end{figure}

由图~\ref{fig:q4_dispatch} 可见，全年储电量始终保持在 $[1200,10800]\,\mathrm{kWh}$ 的允许范围内，并按照跨日连续方式滚动变化。购电量和弃电量随季节变化较为明显，而紧急购电主要集中于少数预测偏差较大的日期，说明日前策略能够满足全年大部分时段的供电需求。

由表~\ref{tab:q4_all}(b) 可见，Q4-3 本文 LA 方案全年总费用为 $14\,001\,199.26$ 元，较不进行日内调整的 N0 减少 $1\,982\,409.26$ 元，降幅约为 $12.40\%$；紧急购电量则由 $466\,576.26$ kWh 降至 $67\,153.26$ kWh，降幅约为 $85.6\%$。与 M0L 相比，LA 仍进一步节省 $825\,631.39$ 元，降幅约为 $5.57\%$。

这表明，在实时电价波动条件下，仅依靠日前计划难以充分应对当天负载和光伏预测误差。利用 $6{:}00$、$12{:}00$ 和 $18{:}00$ 获得的新信息重新调整后续购电计划，可以显著减少高价紧急购电，从而降低全年总费用。

\Needspace{10\baselineskip}
\begin{figure}[!htbp]
  \centering
  \begin{subfigure}[t]{0.48\textwidth}
    \centering
    \includegraphics[width=\linewidth]{q4_cost_structure.png}
    \caption{全年成本结构}
    \label{fig:q4_cost}
  \end{subfigure}
  \hfill
  \begin{subfigure}[t]{0.48\textwidth}
    \centering
    \includegraphics[width=\linewidth]{q4_daily_cost_distribution.png}
    \caption{日总成本分布}
    \label{fig:q4_daily}
  \end{subfigure}
  \caption{Q4-3 三种策略的费用结构与日费用分布}
  \label{fig:q4_cost_daily}
\end{figure}

由图~\ref{fig:q4_cost} 可见，三种策略的费用差异主要来自紧急购电。N0 的紧急购电费用占比较高，而 LA 通过日内滚动调整明显压低了该部分费用，因此全年累计费用始终低于 M0L 和 N0。由图~\ref{fig:q4_daily} 可见，LA 不仅降低了日总费用的整体水平，也明显减少了高费用日期的出现。这说明日内滚动调整的作用不仅体现在降低全年平均费用，还能够减少由预测偏差造成的极端高费用情况。

指定日期的购电量、购电费用与紧急购电量见表~\ref{tab:q4_dates}。

\begin{table}[!htbp]
  \centering
  \caption{波动电价下指定日期的购电结果}
  \label{tab:q4_dates}
  \small
  \begin{tabular}{lrrr rrr}
    \toprule
    日期 & \multicolumn{3}{c}{Q4-2} & \multicolumn{3}{c}{Q4-3} \\
    \cmidrule(lr){2-4}\cmidrule(lr){5-7}
    & 购电量（kWh） & 购电费（元） & 紧急购电量（kWh）
    & 购电量（kWh） & 购电费（元） & 紧急购电量（kWh） \\
    \midrule
    2025-03-20 & $70\,570.86$ & $44\,640.38$ & 0.00
               & $66\,824.83$ & $44\,186.70$ & 165.68 \\
    2025-06-21 & $35\,205.93$ & $17\,564.67$ & 0.00
               & $33\,936.26$ & $16\,638.47$ & 0.00 \\
    2025-09-23 & $68\,848.07$ & $44\,716.63$ & 272.06
               & $67\,439.16$ & $47\,671.18$ & 631.22 \\
    2025-12-21 & $98\,970.30$ & $75\,354.93$ & 0.00
               & $96\,674.19$ & $75\,675.02$ & 291.73 \\
    \bottomrule
  \end{tabular}
\end{table}

由表~\ref{tab:q4_dates} 可见，两类策略的购电需求均表现出明显的季节差异：冬季代表日的全天购电量较高，而夏季代表日明显较低。这主要与全年负载和光伏发电水平的季节变化有关。

同时，个别代表日中 Q4-3 的紧急购电量可能高于 Q4-2，因此不能根据单日结果判断日内调整策略的优劣。结合全年结果可知，LA 将紧急购电量由 N0 的 $466\,576.26$ kWh 降至 $67\,153.26$ kWh，说明其优势主要体现在全年范围内降低预测偏差造成的紧急购电风险。

完整的计划购电量、调整购电量、储能充放电量及紧急购电结果分别保存于 \texttt{result4-2.xlsx} 和 \texttt{result4-3.xlsx}。

\subsection{模型检验}

\subsubsection{电价预测信息检查}

由于问题四的购电决策必须在真实电价尚未完全获得时完成，因此首先检查预测过程中是否使用未来价格信息。对全年 $334$ 个决策日逐日核验后，电价信息泄漏次数为 0，即每日 $0{:}00$ 的预测仅使用此前已经获得的历史电价，$6{:}00$、$12{:}00$ 和 $18{:}00$ 的日内调整也未使用当天真实电价作为优化输入。

该结果表明，全年回测符合题目规定的信息条件，所得费用能够反映实际可执行策略的表现。

\subsubsection{约束一致性检查}

将全年调度结果重新代入储能容量、接口功率及电量平衡等约束进行独立检查。Q4-2 与 Q4-3 的 $334$ 个决策日均满足
\begin{equation}
1200\le E_t\le10800,
\end{equation}
同时满足储能接口功率、电量平衡以及不可向外网售电等运行限制。对各项费用重新按照附件~4 的实时电价计算后，与模型输出结果一致，说明所得调度方案满足运行约束，全年结算结果具有良好的数值一致性。

\subsection{模型对比}

\subsubsection{调度策略对比}

为评价波动电价下本文两类策略的有效性，分别设置与问题二、问题三相同口径的对照方案。Q4-2 将本文净负载分位策略与固定 $q=0.8$ 方案进行比较；Q4-3 则将本文 LA 策略与不进行日内调整的 N0 以及带选择门的 M0L 进行比较。

由表~\ref{tab:q4_all} 可知，Q4-2 本文方案较固定 $q=0.8$ 降低全年费用约 $0.46\%$；Q4-3 中 LA 较 N0 和 M0L 分别降低约 $12.40\%$ 和 $5.57\%$。尤其是 LA 将紧急购电量较 N0 降低约 $85.6\%$，说明在负载、光伏和实时电价均存在变化的情况下，利用日内新增信息调整后续购电计划能够有效降低高价紧急购电造成的费用。

\subsubsection{电价预测影响分析}

为了进一步判断问题四费用增加主要来自实时电价环境本身，还是来自电价预测误差，设置已知当天真实电价的理想对照方案。该方案保持原有调度方法不变，仅将决策时使用的预测电价替换为当天真实电价，并继续按照附件~4 的实时电价结算。

诊断结果如图~\ref{fig:q4_oracle} 所示。

\Needspace{10\baselineskip}
\begin{figure}[!htbp]
  \centering
  \includegraphics[width=0.72\textwidth]{q4_error_decomposition.png}
  \caption{实时电价环境与电价预测误差的费用影响}
  \label{fig:q4_oracle}
\end{figure}

由图~\ref{fig:q4_oracle} 可见，与附件~1 的固定分时电价相比，附件~4 的实时电价环境使理想对照方案的费用提高约 $4.64\%$--$4.76\%$；在此基础上，使用实际可获得的预测信息进行决策仅进一步增加约 $0.58\%$--$0.67\%$ 的费用。因此，问题四相较于固定电价条件下的费用增加主要来源于实时电价环境的变化，而由电价预测误差造成的额外费用相对较小。

这表明本文建立的因果电价预测模型已经能够为购电决策提供较有效的价格信息，同时也说明实时波动电价改变了原有的购电成本环境，使储能利用和购电时机的选择面临更大的价格不确定性。
